% Solutions/hints for ch20 exercises (assembled into Appendix C)

\begin{solution}{exr:ch20-cv-accel}
By \cref{thm:ch20-cv-error} the error at step $h$ is $e_h = \tfrac12 \norm{\acc}(h\,\dt)^2 = \tfrac12 \cdot 2 \cdot (0.2h)^2 = 0.04\,h^2$ metres. Hence $\mathrm{FDE} = 0.04 \cdot 144 = 5.76$~m and $\mathrm{ADE} = \tfrac12\norm{\acc}\dt^2 (H+1)(2H+1)/6 = 0.04 \cdot 13 \cdot 25/6 = 2.17$~m. The error exceeds $0.5$~m when $0.04h^2 > 0.5$, i.e.\ $h^2 > 12.5$, so first at $h = 4$ ($0.8$~s, $e_4 = 0.64$~m). Constant velocity is a safe short-horizon predictor only while $\tfrac12\norm{\acc}(h\,\dt)^2$ stays below the margin; a constant-acceleration predictor would be exact here, but on a turn it overshoots (\cref{fig:ch20-baselines}), so which of the two is safer depends on the manoeuvres that actually occur.
\end{solution}

\begin{solution}{exr:ch20-cv-noise}
Write the observations as $\meas_{T-j} = \pos_T - j\,\dt\,\vel + \vect{\eps}_{T-j}$ for $j = 0, \dots, k$. The least-squares line fit estimates the position at time $T$ and the velocity by linear regression on the times $-j\,\dt$; with $\bar{j} = k/2$ and $S = \sum_j (j - \bar{j})^2 = k(k+1)(k+2)/12$, the fitted velocity is $\hat{\vel} = -\sum_j (j - \bar{j})\meas_{T-j}/(\dt\,S)$ and the fitted position at $T$ is $\hat{\pos}_T = \bar{\meas} + \bar{j}\,\dt\,\hat{\vel}$. Both are unbiased, so the prediction $\hat{\pos}_T + h\,\dt\,\hat{\vel}$ has zero mean error, and its error is a linear combination $\sum_j w_j \vect{\eps}_{T-j}$ with weights $w_j = \frac{1}{k+1} - \frac{(j - \bar{j})(h + \bar{j})}{S}$, giving $\E[e_h^2] = 2\sigma^2 \sum_j w_j^2 = 2\sigma^2\bigl[\frac{1}{k+1} + \frac{(h + \bar{j})^2}{S}\bigr]$. For $k = 7$, $h = 12$: $S = 42$, $\bar{j} = 3.5$, so $\E[e_h^2] = 2\sigma^2 [0.125 + 240.25/42] = 2\sigma^2 \cdot 5.85 = 11.7\sigma^2$, against $2\sigma^2[(19/7)^2 + (12/7)^2] = 20.6\sigma^2$ for the endpoint difference: the line fit reduces the RMS error by a factor $1.33$ ($0.17$ against $0.23$~m at $\sigma = 0.05$~m). For the endpoint difference $2\sigma^2[(1 + h/k)^2 + (h/k)^2]$ decreases monotonically in $k$, so $k = 7$ is best at every $h$ on straight flight. On a turn the velocity rotates during the window, the average velocity over $k$ intervals points $\omega k\,\dt/2$ behind the current heading, and the resulting bias $\approx v\,\omega\,k\,\dt\,h\,\dt/2$ grows with $k$; the validation set of \cref{sec:ch20-example} balanced the two effects at $k = 2$.
\end{solution}

\begin{solution}{exr:ch20-metrics-hand}
The predictions differ from the truth only in $y$, so $e_1 = 0.1$, $e_2 = 0.5$, $e_3 = 1.2$~m, giving $\mathrm{ADE} = (0.1 + 0.5 + 1.2)/3 = 0.6$~m and $\mathrm{FDE} = e_3 = 1.2$~m. The second sample has $e_1 = 0.1$, $e_2 = 0.1$, $e_3 = 0.4$~m, hence an ADE of $0.2$~m and a final error of $0.4$~m. Therefore $\mathrm{minADE}_2 = \min(0.6, 0.2) = 0.2$~m and $\mathrm{minFDE}_2 = \min(1.2, 0.4) = 0.4$~m. Both are minima of a set that contains the first sample's values, so neither can exceed them --- which is why a $\mathrm{minADE}_K$ may never be compared with the ADE of a deterministic predictor.
\end{solution}

\begin{solution}{exr:ch20-attention-complexity}
Projections: $\mat{Q}, \mat{K}, \mat{V}$ each cost $n d^2$ multiply-adds and $nd$ memory. Scores: $\mat{Q}\mat{K}\T$ costs $n^2 d$ and stores $n^2$ numbers; the softmax is $\Theta(n^2)$; the product with $\mat{V}$ is $n^2 d$. Total $\Theta(n d^2 + n^2 d)$ time and $\Theta(n d + n^2)$ memory, and the $n^2$ memory is the usual bottleneck for long sequences. With $n_h$ heads each projection is $n \times d \times d/n_h$ and each score matrix $n^2 d/n_h$; summed over heads the time is unchanged, $\Theta(n d^2 + n^2 d)$, while the score memory becomes $n_h n^2$. For a scene of $N$ agents over $T$ steps as one token set, $n = NT$ gives $\Theta(NT d^2 + N^2T^2 d)$ time and $\Theta(N^2 T^2)$ score memory per layer; per-agent attention over time only costs $N \cdot \Theta(T d^2 + T^2 d)$ and misses the interactions. An LSTM per agent costs $\Theta(T D(M + D)) \approx \Theta(T D^2)$ per agent, $\Theta(N T D^2)$ in total, sequential in $T$ but with no quadratic term. With $d = D = 64$ the joint attention layer costs $64 NT (64 + NT)$ against $64^2 NT \cdot 4$ for the four gate matrices of the LSTM (the factor $4$ accounts for $4D(M + D)$ with $M \ll D$), so attention becomes more expensive when $64 + NT > 256$, i.e.\ when $NT > 192$: about $10$ agents over $20$ steps, or $24$ agents over $8$ steps. Beyond that, attention pays for itself only if it predicts better, and it can be made cheaper by attending over agents at one time step and over time within one agent separately.
\end{solution}

\begin{solution}{exr:ch20-conditional-mean}
Let $\vect{y}^{(1)} = (-2, 4)$ with probability $\pi$ and $\vect{y}^{(2)} = (2, 4)$ otherwise. By \cref{thm:ch20-conditional-mean} the squared-error-optimal end point is the conditional mean $g^\star = \pi\vect{y}^{(1)} + (1 - \pi)\vect{y}^{(2)} = (2 - 4\pi,\, 4)$. Its expected FDE is $\pi\norm{\vect{y}^{(1)} - g^\star} + (1 - \pi)\norm{\vect{y}^{(2)} - g^\star} = \pi \cdot 4(1 - \pi) + (1 - \pi) \cdot 4\pi = 8\pi(1 - \pi)$. For $\pi = \tfrac12$ the prediction is $(0, 4)$, the centre of the junction, with expected FDE $2$~m and a guaranteed miss of $2$~m in every single case; for $\pi = 0.9$ it is $(-1.6, 4)$, expected FDE $0.72$~m, which is $0.4$~m off even when the drone turns left. A Gaussian read-out reports $\mu_x = 2 - 4\pi$ and $\sigma_x^2 = \Var(y_x) = 16\pi(1 - \pi)$, so $\sigma_x = 2$~m for $\pi = \tfrac12$ and $1.2$~m for $\pi = 0.9$: it does not resolve the two modes, but its large $\sigma_x$ announces that something is wrong with a single line. The two-sample predictor that outputs both end points has $\mathrm{minFDE}_2 = 0$, which is why $\mathrm{minADE}_K$ and $\mathrm{minFDE}_K$ are the natural scores for multimodal predictors and why they cannot be compared with the FDE of a deterministic one.
\end{solution}
